\emph{(1).} From $v = \tfrac12(e_{1} + e_{2})$, $w = \tfrac12(e_{1} - e_{2})$: $vw = \tfrac12(1 - e_{1}e_{2})$, $wv = \tfrac12(1 + e_{1}e_{2})$, idempotents with sum $1$ and product $0$ (\claimref{synthesis.set-and-clear-of-a-bit-generate-the-hyperbolic-plane}). Since $e_{2}(e_{1}e_{2})e_{2}^{-1} = -e_{1}e_{2}$, conjugation by $e_{2}$ exchanges $P_{1}$ and $P_{2}$, and $e_{2}P_{1}e_{2} = -P_{2}$. With $U = p + qe_{2}$, $U^{-1} = p - qe_{2}$ ($e_{2}^{2} = -1$, $p^{2} + q^{2} = 1$): $O = U^{-1}P_{1}U = p^{2}P_{1} + q^{2}P_{2} + pq\,[P_{1}, e_{2}]$. As $J = c_{2} + (c_{1} - c_{2})P_{1}$ and $P_{1}$ commutes with $P_{2}$, $[J, O] = (c_{1} - c_{2})pq\,[P_{1}, [P_{1}, e_{2}]]$, and $[P_{1}, e_{2}] = -\tfrac12[e_{1}e_{2}, e_{2}] = e_{1}$, $[P_{1}, e_{1}] = -\tfrac12[e_{1}e_{2}, e_{1}] = e_{2}$, so $[J, O] = (c_{1} - c_{2})pq\,e_{2}$. The source operator is built from the commuting path numbers, as the gravitational source is in Proposition \claimref{field-theory.operator-newtonian-coupling}; before the recombination it is $J$, and the detection, written in the same variables, is $O$.

\emph{(2).} With no free part the field is the retarded solution with source $gJ$ at site $0$ at the times $t'$, which is $gJ\sum_{t'}G(n,t;t')$ by linearity (the overall sign of $g$ is immaterial here); it is a multiple of $J$, so its commutator with $O$ is $g[J,O]\sum_{t'}G$, and two of its values commute. The same computation with any function of the event in place of $\sum G$ gives the last sentence of the clause. The recursion has an exact light cone: a source at $(0,t')$ enters $\phi_{0}(t'+1)$ and reaches $(n,t)$ only if $|n| \le t - t' - 1$. For $t' = 0$ and the event $(7,8)$ this holds with equality, and $(7,8)$ is spacelike to the detection $(0,8)$: the field there is set by the source before the recombination. The count $172$ of $502$ is computed.

\emph{(3).} Both sides satisfy the recursion in $(n,t)$ for $t > t'$, and they agree on the two time slices $t = t'$ and $t = t' + 1$: there $G$ is $0$ and $\delta_{n0}$, and the free field's commutators at equal and adjacent times are $\Delta(n,t';t') = 0$ and $\Delta(n,t'+1;t') = -\delta_{n0}$, the canonical form, which the recursion preserves. For $t \le t'$ the retarded response vanishes. Checked exhaustively on the grid.

\emph{(4).} Let the coupling act as the kicks $K_{t'} = \exp(-I_{4}\,g\,J\,\phi^{0}_{0}(t'))$ in time order. In the Heisenberg picture each kick shifts the field by a multiple $\sigma g J$ of the retarded response (the Yang-Feldman form), so $\phi = \phi^{0} + \sigma g J\sum_{t'}G$ with $\sigma = \pm 1$ determined below, and it conjugates the matter observable, $K^{-1}_{t'}OK_{t'} = O + I_{4}\,g\,[J, O]\,\phi^{0}_{0}(t') + O(g^{2})$, the field commuting with the matter algebra. To first order in $g$,
\[
[\phi_{n}(t), O] = \sigma g\,[J, O]\sum_{t'}G + I_{4}\,g\,[J, O]\sum_{t'}[\phi^{0}_{n}(t), \phi^{0}_{0}(t')] = \sigma g\,[J, O]\sum_{t'}G - g\,[J, O]\sum_{t'}\Delta ,
\]
using $I_{4}^{2} = -1$. Explicitly $K_{t'}^{-1}\phi^{0}_{0}(t'+1)K_{t'} = \phi^{0}_{0}(t'+1) + I_{4}\,g\,J\,[\phi^{0}_{0}(t'), \phi^{0}_{0}(t'+1)] = \phi^{0}_{0}(t'+1) + I_{4}^{2}\,g\,J$, the commutator function being $+1$ there by the canonical form of (3); the shift of $\phi_{0}(t'+1)$ is $-gJ$, so $\sigma = -1$ and the bracket is $-g[J,O]\sum_{t'}(G + \Delta)$. By (3) the terms with $t > t'$ cancel, leaving $\sum_{t' \ge t}\Delta(n,t;t')$, which is nonzero only if $(n,t)$ lies in the past cone of a source event $(0,t')$, $t' \le 5 < t_{d}$; such an event satisfies $|n| \le t' - t < t_{d} - t$ and is timelike to the detection. Hence the vanishing at every spacelike event; the grid statement, the support and the control are computed. The statement is to first order; the working notes observe that it is exact, the kicks being displacement operators, which is not used here.

\emph{(5).} With $\mathrm{Ric}(a) = \sum_{b}e^{b} \cdot \mathcal{R}_{k}[h](e_{b} \wedge a)$ and the Einstein map $\mathrm{Ric} - \tfrac12\,\id\,\tr\mathrm{Ric}$ formed from the curvature map of Theorem \claimref{synthesis.weyl-flat-variations-carry-only-the-stress-trace-and-the-whole-stress-clause-is-the-trace-free-channel}, its rank on the ten-dimensional space of symmetric maps is computed exactly: six at the non-null momenta $(2,0,0,0)$, $(3,1,0,2)$, $(1,2,0,0)$, $(0,1,1,3)$ and four at the null momenta $(1,0,0,1)$, $(5,3,0,4)$. The displacements, of dimension four, lie in the kernel at every momentum (clause (i) of that theorem: they have zero curvature). The rank depends only on the type of the momentum: the map is covariant under the Lorentz group and homogeneous of degree two in $k$, so its rank is the same at $k$, at any Lorentz transform of $k$ and at any nonzero multiple, and every nonzero momentum is of this kind related to a timelike, a spacelike or a null one of the list. So at every non-null momentum the kernel is the displacements, and at every nonzero null one it has two more dimensions. Two inverses of the first-order field equation differ by a solution without source. The free field of (4) obeys $[\phi^{0}, \phi^{0}] = I_{4}\Delta$, a canonical commutation relation, which by Proposition \claimref{field-theory.no-canonical-commutation-relations-tower} no floor of the tower contains.

\emph{Reading.} Theorem \claimref{synthesis.locality-of-the-attached-free-fields-fixes-their-statistics} concerns the bracket of an attached free field with itself at spacelike separation. In (2) every value of the field is a multiple of $J$, so that bracket is zero, while $[\phi, O] \neq 0$ at spacelike events: the two conditions are different, and the second is not stored. The stored coupling of Proposition \claimref{field-theory.operator-newtonian-coupling} is static; taking the field at every event to be a function of the matter's numbers is the extension made in the Reading. For a field $gJf$ with any function $f$ of the event the commutator with $O$ is $g[J,O]f$, nonzero wherever $f$ is, which is the remark on the instantaneous field. The rest of the Reading is the standard argument and is labelled argued.
